\section{从固定配比转向风险预算优先}

此前10\%／15\%／40\%／25\%／10\%参考向量包含40\%红利权益ETF和25\%公募REIT，已知风险资产层的规划压力小计为-25.00\%，高于-18\%风险行动线。该算术说明预设类别权重没有通过风险预算；它不应被修补为另一套未经验证的固定配比，更不能据此永久排除直接股票。该向量现仅保留为“被否决的历史研究设计”，用于说明为什么未来组合必须从获批标的的边际风险贡献出发。

\begin{exhibitbox}[历史参考向量的共同压力：用于否决固定配比，不是新组合的损失预测]
\small
\begin{tabularx}{\exhibitboxwidth}{L{3.1cm}P P P}
\toprule
\textbf{历史研究层} & \textbf{历史参考权重} & \textbf{规划压力输入} & \textbf{已知加权贡献} \\
\midrule
受保障流动性 & 10\% & 0\% & 0.00\% \\
五年期国债ETF & 15\% & -5\% & -0.75\% \\
红利权益ETF & 40\% & -45\% & -18.00\% \\
公募REITs & 25\% & -25\% & -6.25\% \\
\midrule
\textbf{已知资产层小计} & \textbf{90\%} &  & \textbf{-25.00\%} \\
未定义防御分散器 & 10\% & 无获批输入 & 不得纳入通过计算 \\
\bottomrule
\end{tabularx}
\end{exhibitbox}

该表不是历史回撤、VaR、预期损失、概率或最大损失。它只说明旧的类别向量不具备激活资格。下一次组合评审需由T3批准可投池的具体标的生成新的压力、相关性、流动性和总回报归因模型；若风险预算不批准，正确结果仍是$q_{portfolio}=0$，而不是把资金转向任何“更高收益”替代品。

\section{开放研究宇宙下的硬性风险控制}

\begin{exhibitbox}[所有资产格式共用的组合约束]
\small
\begin{tabularx}{\exhibitboxwidth}{L{3.15cm}P P}
\toprule
\textbf{控制} & \textbf{适用逻辑} & \textbf{失败结果} \\
\midrule
单项研究完整性 & 每项必须有结构适配的R0--R4档案、当前估值/参考值、现金机制、风险、催化和证伪。 & 不能进入T3。 \\
共同风险映射 & 直接股票、ETF与REIT按发行人集团、行业、因子、利率、客户/租户/承购方、融资和地域/监管共同驱动聚合。 & 未能映射的候选$q=0$。 \\
边际风险贡献 & 新候选须改善或至少符合获批压力损失、集中度、流动性和容量预算；单独高股息不是充分条件。 & 不进入T4。 \\
分配与现金流质量 & 核对分配削减、股利／经营现金流覆盖、返还性成分、税费和成本，以判断经济质量与总回报归因。 & 不得把资产处置所得、杠杆、低信用或强制高收益配置伪装为回报改善。 \\
流动性与实施 & 审核A股自由流通量、涨跌停、停复牌、价差、深度、参与率、最小单位、税费与结算。 & 无可执行性即$q=0$。 \\
资产类别特别授权 & 杠杆、反向/杠杆产品、期货/期权、私募或显著低流动性资产需要独立书面授权。 & 不得由开放宇宙自动纳入。 \\
\bottomrule
\end{tabularx}
\end{exhibitbox}

\section{观察回撤的强制风险动作}

观察回撤仅用带日期、外部现金流调整后的账户市值和高水位计算：
\[
DD_t=\frac{MV_t^{\mathrm{flow\ adjusted}}}{HWM_t^{\mathrm{flow\ adjusted}}}-1.
\]

\begin{exhibitbox}[风险行动线不是被动“复核”]
\small
\begin{tabularx}{\exhibitboxwidth}{L{3.2cm}P}
\toprule
\textbf{触发} & \textbf{强制治理结果} \\
\midrule
$DD_t>-12\%$且其他闸门均通过 & 仅可讨论已在T3且有效$q=1$的候选；不意味着自动下单。 \\
$DD_t\leq-12\%$ & \textbf{谨慎锁定}：所有风险新增暂停；若投资者另行考虑配置，须先重新承保。 \\
$DD_t\leq-18\%$ & \textbf{风险停止}：所有新的市场风险配置均禁止；取消未执行模型批次，形成事件与去风险方案后方可讨论再批准。 \\
达到获批压力容量的90\%（若未来有例外） & 风险停止继续生效，重做持仓、流动性和资金用途风险预算；无机械重入。 \\
\bottomrule
\end{tabularx}
\end{exhibitbox}

这些是新资金的治理权限，不是自动卖出指令，也不保证亏损在阈值停止。悬停、跌停、跳空和低流动性情形仍可能使实际损失超过行动线。

\section{产品结构专项控制不等于类别配额}

公募REIT在进入T4时，仍须满足至少五个独立项目、三个以上管理人/发起人集团和三个以上运营现金流行业；单项目、单一管理人/发起人和共同风险簇必须受组合上限约束。红利ETF与直接股票则需要发行人集团、行业和因子集中度映射。直接股票不因个股数量不同而自动分散，三只租赁住房REIT也不因代码不同而天然独立。

未来可由风险预算为单一发行人、单一基金、单一REIT项目、行业/因子/现金流风险簇分别设定上限，但本案没有获批的候选组合或损失容量，因此不虚构数值上限。唯一有效的当前控制是：全部市场候选保持$q=0$，直至独立风险复核完成。
